\section{Uniqueness of the folding-map family under protocol locality (closure)}
\label{app:fold_family_uniqueness}

This appendix closes (OP2) by recording a minimal ``protocol-local/implementable'' notion of a folding-map \emph{family} and proving that, under this notion, the folding bridge is uniquely forced to be the Zeckendorf-truncation family $\mathrm{Fold}_m$ (Definition~\eqref{eq:foldm_def}).
The key point is that OP2 is not about whether alternative deterministic maps $\{0,\dots,2^m-1\}\to X_m$ exist at a \emph{fixed} window length (they do, and Appendix~\ref{app:fold_family_sensitivity} audits a bounded family at $m=6$), but about whether a \emph{principled} choice is forced once one commits to a locality/implementability contract for how the dyadic register is refined.

\begin{center}
\fbox{\begin{minipage}{0.96\linewidth}
\small
\textbf{Audit note (status).}
The result below is a closure \emph{within} an explicit contract:
\emph{value consistency} (Definition~\ref{def:value_consistency_m}) together with a \emph{no-backreaction refinement rule} (Definition~\ref{def:protocol_local_fold_family}).
This does not claim invariance under arbitrary dyadic$\to X_m$ repairs; rather, it shows that once one requires a protocol-local family compatible with window uplift, the only possibility is the Zeckendorf-truncation family already used in the folding core.
\end{minipage}}
\end{center}

\subsection{Protocol-local folding families}
\label{subsec:protocol_local_fold_family}

We write $X_m\subset\{0,1\}^m$ for the admissible language (no adjacent ones) on the golden branch, and we use the prefix projection (Appendix~\ref{app:functorial_refinement})
\[
\pi_{m\to k}:X_m\to X_k,\qquad
\pi_{m\to k}(w_1\cdots w_m):=w_1\cdots w_k
\qquad (m\ge k\ge 1).
\]
We also use the natural set inclusion $\{0,\dots,2^m-1\}\subset\{0,\dots,2^{m+1}-1\}$ to regard an $m$-bit dyadic index $N$ as the same integer when a longer register is available.

\begin{definition}[Protocol-local folding-map family]\label{def:protocol_local_fold_family}
A \emph{protocol-local folding-map family} is a collection of deterministic maps
\[
F_m:\{0,\dots,2^m-1\}\to X_m,\qquad m\ge 1,
\]
such that:
\begin{enumerate}
\item \textbf{Value consistency.} For every $m\ge 1$, the map $F_m$ is value-consistent in the sense of Definition~\ref{def:value_consistency_m}, i.e.
\[
F_m\!\bigl(V_m(w)\bigr)=w\qquad\text{for all }w\in X_m,
\]
where $V_m(w)=\sum_{k=1}^m w_k F_{k+1}$ is the Zeckendorf value used in Proposition~\ref{prop:foldm_surjective}.
\item \textbf{Refinement locality (no backreaction under dyadic uplift).}
For every $m\ge 1$ and every $N\in\{0,\dots,2^m-1\}$, the $m$-digit output at window length $m$ is the prefix of the $(m+1)$-digit output at window length $m+1$:
\begin{equation}\label{eq:fold_family_refinement_locality}
\pi_{m+1\to m}\!\bigl(F_{m+1}(N)\bigr)=F_m(N).
\end{equation}
\end{enumerate}
\end{definition}

\paragraph{Interpretation.}
Equation \eqref{eq:fold_family_refinement_locality} is a minimal implementability constraint: increasing the dyadic register budget by one bit is allowed to \emph{refine} the stable word by adding one more admissible digit, but it is not allowed to rewrite the previously reported prefix.
Equivalently, the outputs form a projective system under the truncation maps $\pi_{m+1\to m}$, in the same sense that the label lift $\mathcal{L}_m=\mathcal{L}_{\mathrm{SM}}\circ\pi_{m\to 6}$ is functorial under prefix refinement (Appendix~\ref{app:functorial_refinement}).

\subsection{Uniqueness theorem}
\label{subsec:fold_family_uniqueness_theorem}

\begin{theorem}[Uniqueness of $\mathrm{Fold}_m$ under protocol locality]\label{thm:fold_family_uniqueness}
Let $(F_m)_{m\ge 1}$ be a protocol-local folding-map family in the sense of Definition~\ref{def:protocol_local_fold_family}.
Then for every $m\ge 1$ and every $N\in\{0,\dots,2^m-1\}$ one has
\[
F_m(N)=\mathrm{Fold}_m(N),
\]
where $\mathrm{Fold}_m$ is the Zeckendorf-truncation folding map defined in \eqref{eq:foldm_def}.
In particular, the folding-map family is unique under the stated contract.
\end{theorem}
\begin{proof}
Fix $m\ge 1$ and $N\in\{0,\dots,2^m-1\}$.
Choose an integer $M\ge m$ such that
\begin{equation}\label{eq:choose_M_for_value_consistency}
N<F_{M+2}.
\end{equation}
Such an $M$ exists because $F_{k}\to\infty$ as $k\to\infty$.

Let $w:=\mathrm{Fold}_M(N)\in X_M$ be the length-$M$ Zeckendorf digit prefix of $N$ (Definition~\eqref{eq:foldm_def}).
Condition \eqref{eq:choose_M_for_value_consistency} implies that the Zeckendorf expansion of $N$ uses no Fibonacci weights beyond $F_{M+1}$, hence the digits beyond position $M$ are zero and the length-$M$ digit string $w$ is the full Zeckendorf digit string of $N$ padded by zeros.
Equivalently,
\begin{equation}\label{eq:VM_of_foldM_is_N}
V_M(w)=N.
\end{equation}

By value consistency (Definition~\ref{def:protocol_local_fold_family}(1) and Definition~\ref{def:value_consistency_m}),
\[
F_M(N)=F_M\!\bigl(V_M(w)\bigr)=w=\mathrm{Fold}_M(N).
\]
Applying refinement locality \eqref{eq:fold_family_refinement_locality} repeatedly (from $M$ down to $m$) yields
\[
F_m(N)=\pi_{M\to m}\!\bigl(F_M(N)\bigr)=\pi_{M\to m}\!\bigl(\mathrm{Fold}_M(N)\bigr)=\mathrm{Fold}_m(N),
\]
since truncating the first $M$ Zeckendorf digits to length $m$ is exactly the definition of $\mathrm{Fold}_m$.
\end{proof}

\subsection{Relation to the bounded counterfactual audit at \texorpdfstring{$m=6$}{m=6}}
\label{subsec:fold_family_uniqueness_relation_m6}

Appendix~\ref{app:fold_family_sensitivity} records a bounded family of deterministic counterfactual maps at $m=6$.
Within that fixed-$m$ family, value consistency already selects FoldZ uniquely (Proposition~\ref{prop:value_consistency_selects_foldz}), and within the shifted Zeckendorf-window family it forbids all nonzero shifts (Proposition~\ref{prop:value_consistency_forbids_shift}).
Theorem~\ref{thm:fold_family_uniqueness} upgrades this to a global statement: once one additionally requires refinement locality across window lengths (the minimal implementability contract), the only possible folding-map family is $\mathrm{Fold}_m$.

\subsection{Scope (what is not claimed)}
\label{subsec:fold_family_uniqueness_scope}

The closure here is \emph{not} an invariance claim under arbitrary changes of the dyadic$\to X_m$ bridge.
It is a principled selection theorem: within the stated protocol-local family class, the bridge is uniquely forced and introduces no residual freedom.
If one drops refinement locality or value consistency, then many alternative families exist (one can modify the map on dyadic indices outside $\{0,\dots,F_{m+2}-1\}$ at fixed $m$, or allow backreaction under uplift), and OP2 would remain genuinely open.


